\chapter{Why States Appeared}
\section{A Ten-Yuan Piece of Paper}
Pick something up: a ten-yuan banknote.

Inspect its decent paper, landscape, patterns, national emblem, and "People's Bank of China". Manufacturing probably costs no more than a few tenths of a yuan. Yet any shop will exchange it for noodles, a soft drink, or two bus rides.

Strange: noodles are edible and filling; the note is paper. Why?

Neither paper's value nor attractive printing explains it. Everyone believes it useful because something stands behind it: the state. The state calls it money, accepts it for taxes, prohibits refusal. Collapse that backing and paper reappears. History repeatedly turns banknotes worthless. In 1940s China, government Gold Yuan notes depreciated within less than a year until sackfuls bought rice, then rice merchants refused them. Same fine printing, no guarantee, just paper.

Enlarge the question. Where did the state come from, this guarantor of paper and mutual confidence among millions of strangers?

It may seem always present: histories overflow with dynasties, kings, emperors. Extend the timescale and be startled. Homo sapiens has lived two or three hundred thousand years; earliest states appeared around five thousand years ago. For roughly ninety-eight per cent of human existence there were no kings, written laws, tax offices, police, identity cards. Then, between roughly five and three thousand years ago, states independently arose in Mesopotamia, Egypt, the Indus region, China, Mesoamerica, and the Andes like separate mushrooms after rain.

Why, after over two hundred thousand stateless years, this historically sudden convergence? Did people want states, or did states capture people? That is our hard question.

\section{A Day at the Ancient Warehouse Door}
Enter a scene.

Time: around 3200 BCE. Place: Uruk beside the lower Euphrates in today's Iraq, among the world's largest settlements, perhaps tens of thousands strong.

Early sun already warms earthy river smells. A queue forms outside the central temple store. Men drive barley-laden donkeys or carry beer jars. Donkeys snort; drivers mutter curses. Progress is slow because seated scribes work over every delivery.

No paper, whose invention lies millennia ahead, but palm-sized wet tablets. Pointed reeds press small signs for beer, barley, owner, date. Drying or baking hard fixes the record. Thousands of excavated tablets largely record receipts, disbursements, balances. Earliest writing contains neither poems nor history, but accounts.

One recurring sign on the oldest accounts is sometimes read as Kushim. It appears on dozens of tablets handling quantities of beer and grain. Researchers dispute personal name versus office, perhaps warehouse supervisor. Either way, meaningful: among history's first named people may stand neither hero nor king, but an accountant.

At the store's other end, afternoon brings another queue: temple artisans, porters, brewers collecting rations. Scribes press fresh marks for barley issued. Every movement in or out leaves a record.

Notice the whole machine: rule makers fixing dues and allowances, recordkeepers, guards preventing grain's disappearance, temple priests and urban leaders issuing commands. This is a state's embryo. Its first component is ledger, not sword.

\section{Not a Question of the First King}
Correct an easy misdirection.

We ask not who first became king, but something harder and more intimate: for whom was this machine built?

Two irreconcilable pictures stand before you.

First: saviour. Tens of thousands need irrigation allocation, protection against upstream monopolisation, injury compensation, walls, flood and raider defence. Without coordination, city life seems impossible. Thomas Hobbes's seventeenth-century Leviathan calls life without common power solitary, poor, nasty, brutish, and short. People surrender some freedom and force through a great contract, exchanging it for collective state protection. Taxes become premiums, officials property managers, laws contract terms.

Second: robber. Why must farmers surrender grain? Every tablet impression records your debt to authority. Failure reveals whips, forced labour, military service. Historian Charles Tilly's twentieth-century formulation, war makes states and states make war, likens government to protection rackets: manufacture or exaggerate threats, then charge for protection. Uruk's store becomes a till, its tablets IOUs.

Both have evidence and claim common sense. Their gap holds our question:

\textbf{Is the state a tool people invented to survive together, or a cage the few built to control the many?}

Before choosing, hear both those feeding and serving the machine and those operating it.

\section{Tax Collectors and Fugitives}
\textbf{First give the state its full case.} Calling it oppression wins easy applause today; precisely that ease makes examining the strongest opposing reasons essential.

Imagine being Uruk's store supervisor or a later dynasty's county magistrate. Your defence: collected barley does not simply enter my pocket; corruption, where present, distorts rather than defines the institution. Emergency grain saves last year's tax critics. Who coordinates thousands repairing a breached embankment? Neither one household nor ten can do it; without labour-mobilising authority the plain drowns. Who protects desert caravans, suppresses bandits, maintains stations? Without public adjudication, one injury can exhaust two villages in hereditary vendettas. Courts let judgement end grievance. Your coercion is my cost of order; your taxes fund roads, troops, disaster stores. Without us, you might not survive to complain.

Strong reasoning. Serious political-science textbooks still place security and public goods among the state's primary justifications.

\textbf{Now give the other side its full case.}

Imagine paying grain from riverside fields instead. Good harvests lose substantial shares; bad harvests still owe dues. Borrowing compounds interest until you mortgage yourself into debt bondage. A son spends months building walls while family fields lack labour. Your calculations never balance dues with returned benefits. Temple gold and silver accumulate; your few pottery bowls stay the same.

A more complete answer is flight. Political anthropologist James Scott, studying Southeast Asia in the later twentieth century, describes in The Art of Not Being Governed upland peoples of what he calls Zomia, many descended from farmers deliberately fleeing lowland states. They move uphill, adopt crops not needing central irrigation, even deliberately retain oral traditions because writing serves taxes and conscription. They are not primitives awaiting state invention, but informed people comparing prices and voting with their feet.

Beware equating statelessness with stupidity, laziness, backwardness. Stateless societies have order, arbitration, affluence of their own. Anthropological estimates suggest some hunter-gatherers spend twenty or thirty weekly hours obtaining food, with considerable regional variation, leaving conversation, song, dance, rest. States are no graduation from savagery into civilisation; many experience them as an unsolicited bill.

Both voices heard, notice their conflict concerns \textbf{readings} of shared facts, not denying that grain is distributed and taxes collected. Progress requires cleaning several large words.

\section[Thinking Bridge: Separate Four Terms]{A Thinking Bridge: Take Four Big Words Apart}
State, city, monarchy, empire often blur in news, textbooks, conversation. Confusion makes disputes talk past one another from the start. Our simple tool: \textbf{take big terms apart before debating them.}

\begin{itemize}
\item \textbf{City: a density of settlement.} Many people live together through exchange and specialisation rather than directly obtaining food from land. It describes density and livelihood.
\item \textbf{State: an administrative machine.} Its core is not population size but specialised tax, archival, enforcement, military institutions and territorial authority. It describes who decides and through which machinery.
\item \textbf{Monarchy: an arrangement of supreme leadership.} Highest power belongs to one person and family, usually hereditary and sacred. It is one configuration of a state, not the state itself.
\item \textbf{Empire: a territorial structure.} A core state brings diverse peoples and lands under one rule, often unequal, extracting from peripheries. It describes how many different populations the machine encompasses.
\end{itemize}
Recombine them and none automatically entails another:

\begin{tabularx}{\linewidth}{llX}
\toprule
Combination & Exists? & Example \\
\midrule
Cities, no evident monarchy & Yes & Harappa and Mohenjo-daro in the Indus region \\
Monarchy, no cities & Yes & Seasonally mobile Eurasian nomad courts \\
State, no monarchy & Yes & Classical Athens under its citizen assembly \\
State, no writing & Yes & Inca knot-record administration \\
Empire & = State + diverse ruled peoples & Assyria, Persia, Rome, unified dynasties \\
\bottomrule
\end{tabularx}
Inspect the most counterintuitive cases.

Harappa and Mohenjo-daro yield planned cities with straight streets, standard bricks, almost luxurious drains, yet no recognised palaces, royal tombs, or monuments celebrating war and kingship. Organised urban life seems present without identifiable kings. Honestly, undeciphered Indus writing leaves vast ignorance; absence of kings is only a hypothesis from current evidence.

Eurasian nomad khans could command whole steppes while seasonal courts travelled on carts. Monarchy here need not connect to a city.

Classical Athens debated affairs in the assembly and chose many officials by lot. Today's republics lack kings too; state machines work with another answer to leadership.

Most misleading is no writing, no state. The Inca governed over ten million across thousands of mountain kilometres, building roads, mobilising labour, allocating goods, without writing in our sense. Quipu used colours, knot positions, strand numbers to record quantities and affairs under specialist keepers. States need information storage and transmission; writing is one method.

Empires add a layer: one machine over different peoples. Assyria, Persia, Rome, and unified Chinese dynasties from Qin and Han illustrate it.

Why bother? Confused words create confused accounts. Debating unification without separating monarchy and state equates rejecting emperors with rejecting public order. Blurring state and empire equates national belonging with support for domination of others.

Remember: at a big word, ask whether it concerns density, leadership, machinery, or territory.

\section[Read an Ancient Law Code]{Read a Thirty-Seven-Century-Old Law Code}
Now primary material.

In late 1901 and early 1902, French excavators at Susa in today's Iran found a black basalt stele over two metres tall. Its relief shows a standing figure before a seated god, usually identified as Hammurabi receiving a sceptre from sun and justice god Shamash. Below lie over 280 cuneiform laws. This eighteenth-century BCE Code of Hammurabi, whose reign ran approximately 1792--1750 BCE, now stands in the Louvre. Earlier written laws include Ur-Nammu's code under the Third Dynasty of Ur; Hammurabi's is among the most complete and famous, not the first.

Its preface claims divine selection to make justice appear, destroy wrongdoing, and prevent strong oppressing weak. Remember that self-definition; we will inspect it again.

The body uses if--then provisions. Try three paraphrases:

\begin{quote}
If someone neglects an embankment and its breach floods a neighbour's field, he must compensate the neighbour's grain. (Around section 53.)
\end{quote}
\begin{quote}
If someone destroys another free man's eye, his own eye shall be destroyed. (Section 196.)
\end{quote}
\begin{quote}
If a builder constructs an unsound house whose collapse kills its owner, the builder shall die. (Section 229.)
\end{quote}
Pause over three major points.

\textbf{First, law transforms custom into public rules.} People already repaid debts, compensated injury, and answered for failed dykes, but rules lived in elders' memory, village custom, priests' speech. Flexibility carried weaknesses: differing memories, changing village norms, powerful people declaring rules without weaker parties' means to check. Stone makes rules \textbf{public property}, in principle open to inspection, challenge, citation. From what people say to what is written comes a quiet political earthquake. The strong do not automatically improve, but the weak gain a weapon: pointing to a provision.

\textbf{Second, the code engraves existing inequality too.} Notice free man's status. Damage to one such person's eye costs an eye; damage to a lower-status free mushkenu costs silver; enslaved people resemble owners' property, with compensation paid to owners. Justice assigns unequal prices. Public writing fixes \textbf{contemporary} rules, fair and unfair alike, with stone lending greater authority to both.

\textbf{Third, ideal order is not daily practice.} A publicly or temple-displayed stele crowns itself with divine authority and royal proclamations of justice. Many researchers therefore read \textbf{a declaration of achievements and legitimate kingship}, not judges' working manual. Excavated court records show decisions according to circumstances and local practice, rarely verbatim code citations. We read the state Hammurabi \textbf{wanted}, his blueprint of protected weakness. Real Babylon still held debt slaves, deserters, corrupt officials, clever litigators. Reading codes as surveillance footage is a common historical trap.

\section{Five Explanations, Five Keys}
Why states arose has five principal explanatory keys. Every following claim concerns why, not factual occurrence or moral evaluation.

\textbf{Explanation 1: Violence and war.} Armed conquerors discover annual extraction beats daily robbery. Settled bandits replace mobile ones, crown themselves, become dynasties. Tilly's European account makes sustained war develop taxes, recruitment, administration; losers unable to compete disappear. Sharp for dynastic turnover and expansion, it poorly explains early apparently unconquered temple cities or popular acceptance beyond rulers' capacity to dominate.

\textbf{Explanation 2: Irrigation.} An influential twentieth-century hypothesis associated with Karl Wittfogel links river-plain agriculture to large waterworks requiring supra-village authority: water breeds bureaucracy, then state. The code's grain-compensation rule seems supportive. Yet here is our \textbf{easily misread counterexample}: many major Mesopotamian and Egyptian irrigation schemes follow state formation, not precede it. Early works were locally manageable. Water needs coordination, but projects-before-authority often reverses chronology. The smoother the explanation sounds, the harder to prod its timeline.

\textbf{Explanation 3: Trade.} Southern Mesopotamia lacks good stone, metal ores, timber, but has exceptionally fertile fields. Temples and tools require exchanging grain for distant materials. Long-distance trade needs measures, contracts, guards, arbitration, permanent institutions growing into states. Trade nodes elsewhere also develop cities and kingship early. Yet trade itself presupposes order; it may fertilise state growth rather than seed it.

\textbf{Explanation 4: Religion and legitimacy.} Uruk's stores, records, workers belong to \textbf{temples}; Hammurabi receives authority from the sun god; Shang kings govern through ancestral and divine divination; Zhou explain dynastic change through Heaven's Mandate, withdrawn from immoral rulers and granted to virtuous ones. Religion is an early \textbf{licensing authority}, answering why this ruler, a question force cannot settle alone. It supplies obedience's missing explanation, but leaves timing vague and may merely package successful violence afterward.

\textbf{Explanation 5: Population and collective coordination.} Farming increases numbers crowded onto shared plains, disputing water, land, boundaries beyond kin mediation and reputation's capacity. Arbitration requires full-time specialists, support, grain collection; scribes, stores, standing forces follow. This makes states practical responses without presumed goodness or evil. Yet it renders them too harmless: how does a salaried mediator become hereditary, divine, irremovable king? Its gentle explanatory slope cannot cross that gap.

Each key opens some doors, none all. Most researchers now favour regional recipes: Mesopotamian temples and trade, Egyptian water and kingship, steppe conquest and trade routes, Mesoamerican ritual centres and population pressure. One universal cause assumes one worldwide lock. Nor were states inevitable: humanity spent ninety-eight per cent of its time elsewhere; Zomian peoples could leave, return, leave again. States are historical products, not human destiny.

\section[Further Challenge: Legitimate Force]{A Deeper Challenge: The Legitimate Monopoly of Violence}
Our strongest theoretical tool comes from German sociologist Max Weber. In his famous 1919 Munich lecture Politics as a Vocation, he defines the state, in paraphrase, as \textbf{a human community successfully claiming monopoly over legitimate physical force within a territory}.

Every word deserves chewing over.

\textbf{First, violence.} Whatever its packaging, coercion is the state's final card. Violation brings summons, judgement, fines, seizure, bailiffs, prison. Civilised procedures can stretch the chain enormously, but its endpoint remains someone authorised to use force. Missing homework may summon parents; state-level refusal cannot continue indefinitely. This does not mean constant blows. Successful states rarely strike because everyone knows the final card.

\textbf{Next, monopoly.} Confining a person becomes lawful imprisonment for the state, unlawful detention for you. Taking earnings becomes tax for the state, robbery for others. It claims \textbf{exclusive authority}, including authorised delegates, to use force legitimately. Other violence becomes unlawful, including once-accepted revenge and debt collection. This monopoly forcibly interrupts hereditary vengeance.

\textbf{Finally, legitimacy, the subtlest term.} Anyone may want monopoly; why accept this claimant? Weber distinguishes \textbf{tradition}, the family always ruled; \textbf{charisma}, an extraordinary leader worth following; and \textbf{legal-rational authority}, selection and action under rules. Earlier threads converge. Divine Hammurabi, Zhou's Mandate, modern constitutions, elections, inaugurations all work for legitimacy. Force monopoly supplies the skeleton; legitimacy the blood that keeps it alive.

This framework connects Uruk's taxing stores, administrative tablets, Hammurabi's public claim, Tilly's interstate competition, and Zomian freedom beyond the monopoly's reach.

Its chilling implication deserves honesty: \textbf{a functioning state and a protection racket are structural relatives}. Both charge, protect, exclude competitors' force. Tilly nearly calls states the most successful protection businesses. Where is the difference? Public fees and accounts? Protection delivered? Replaceable collectors? Channels for disagreement and exit? Not whether payment exists, but these answers. The theory defends neither state nor gang; it hands you a ruler: \textbf{measure each authority's distance from public service and proximity to protection racketeering}.

Remember limits: monopoly often fails; civil war pits two state claimants against one monopoly. Legitimacy erodes; widespread rejection leaves a skeleton vulnerable to collapse. Weber describes, not praises.

\section{You Live Daily Inside the State's Body}
This ancient question operates on you now.

Your schoolward road is publicly built or planned, tendered, inspected. Your school is probably public, textbooks approved, teachers publicly paid. Compulsory education is among the state's greatest public goods, so familiar you need reminders to see it. Electricity, water, public hospitals, police calls all connect to the same machine.

You feed it too, perhaps unnoticed. Drink and trainer prices include taxes; parents' wages lose a share before receipt. Later come social insurance, income tax, military registration. Protection has never been free, its unchanged five-millennia charter.

School offers miniature state history. A new class begins with habits: whoever is available cleans the board, whoever is loud keeps order. Weeks later habits fail: why always the same cleaner? Written wall rules turn \textbf{custom into public regulation}, your class's law code. Class leader, subject representatives, duty roster bring bureaucracy; disputed basketball decisions go to the teacher's court; class funds, spending, publication bring taxes and budget struggles. Every problem resembles Uruk's store entrance: collect public money, divide public labour, adjudicate disputes, justify authority.

Measure something new: \textbf{online platforms}. Territory, servers and users; law, lengthy agreements and community rules; enforcement, deletion, throttling, account bans nearly digital expulsion; taxes, commissions; even currencies of coins and beans. States? Weber's ruler finds no legitimate physical-force monopoly: banning cannot arrest. Their legitimacy rests on a consent button, yet control over hundreds of millions' expression, trade, reputation exceeds many state departments. What twenty-first-century polity is a platform? The world lacks a good answer, but your analytical ruler may matter more.

Practise our distinctions again. Public education and security are facts. State-as-only-domination and state-as-only-public-service are rival explanations. Your future tax, service, and authority choices concern evaluation and action. This chapter fixes the first layer and opens the second; you gradually develop answers for the third.

\section[Your Turn: What Price Is Fair?]{Your Turn: What Is a Fair Price for This Bargain?}
Return to the ten-yuan note.

It buys noodles because you, shopkeeper, printers, banks, courts, millions of strangers share and broadly recognise one machine. To sustain it, you surrender earnings, some freedoms, including once-accepted private revenge, obedience, and sometimes in some countries years of youth or life. It returns roads, schools, order, disaster grain, and ordinary days without fear of routine robbery.

Five millennia of bargaining weigh freedom against security.

Now answer: \textbf{is what the machine takes fairly priced against what it gives?} Sharper: tomorrow an exit button offers no taxes, service, protection, or jurisdiction, moving beyond its reach like Zomian highlanders. Press it or not?

Before giving reasons, count your counters.

For the machine: Uruk's stores opened during scarcity; breached dykes really incurred compensation; absent common power, vendettas can persist for generations. Your peaceful reading rather than debt labour or conscription flight itself indicates successful machinery.

For fugitives: engraved inequality proclaims protection of the weak; collectors always set protection's price; informed highlanders voted with their feet. Majority recognition can mean legitimacy or merely entrenched habit. How different are you from residents paying a gang monthly?

Use our tools: distinguish facts, interpretation, evaluation; separate city, monarchy, state, empire; apply Weber's ruler to every authority, including the obvious assumptions in your answer.

No standard answer. One last point: nearly all Zomian descendants today have descended, carrying identity cards and ten-yuan-type notes. Were they wrong? Did the price change? Or was there never a genuine button?

What are your reasons?
